\documentclass[11pt]{article}

\usepackage{amsmath, amssymb, amsthm}
\usepackage[margin=1in]{geometry}

\newtheorem{theorem}{Theorem}

\DeclareMathOperator{\disc}{disc}

\title{Discrepancy of the Complete Set System}
\author{}
\date{}

\begin{document}

\maketitle

\begin{theorem}
Let \(X\) be a finite set with \(|X| = n\), and let
\(\mathcal{F} = \mathcal{P}(X)\) be the complete set system on \(X\),
i.e.\ the collection of all subsets of \(X\). Then the discrepancy of
\(\mathcal{F}\) is
\[
\disc(\mathcal{P}(X)) = \left\lceil \frac{n}{2} \right\rceil.
\]
\end{theorem}

\begin{proof}
Recall that the discrepancy of a set system \(\mathcal{F}\) on a
ground set \(X\) is defined by
\[
\disc(\mathcal{F})
= \min_{\chi \colon X \to \{\pm 1\}}
  \max_{S \in \mathcal{F}}
  \left| \sum_{x \in S} \chi(x) \right|,
\]
where the minimum is taken over all two-colorings of \(X\), and the
maximum over all sets in the system.

We prove the two requisite inequalities separately.

\medskip
\noindent\textbf{Lower bound.}
Fix an arbitrary coloring \(\chi \colon X \to \{\pm 1\}\), and put
\[
P = \{x \in X : \chi(x) = +1\},
\qquad
N = \{x \in X : \chi(x) = -1\}.
\]
Let \(p = |P|\) and \(q = |N|\), so \(p + q = n\). Since
\(\mathcal{P}(X)\) contains every subset of \(X\), it contains \(P\)
and \(N\) in particular. Therefore
\[
\max_{S \in \mathcal{P}(X)}
\left| \sum_{x \in S} \chi(x) \right|
\ge
\max\left\{
  \left| \sum_{x \in P} \chi(x) \right|,
  \left| \sum_{x \in N} \chi(x) \right|
\right\}
= \max\{p, q\}.
\]
As \(p + q = n\), we have \(\max\{p, q\} \ge \lceil n/2 \rceil\).
Since the coloring \(\chi\) was arbitrary, it follows that
\[
\disc(\mathcal{P}(X)) \ge \left\lceil \frac{n}{2} \right\rceil.
\]

\medskip
\noindent\textbf{Upper bound.}
It remains to exhibit a coloring whose maximum discrepancy over all
subsets is at most \(\lceil n/2 \rceil\). Choose a coloring \(\chi\)
for which
\[
|P| = \left\lceil \frac{n}{2} \right\rceil,
\qquad
|N| = \left\lfloor \frac{n}{2} \right\rfloor.
\]
Such a coloring evidently exists. Now let \(S \subseteq X\) be
arbitrary, and write \(S = S_+ \cup S_-\), where
\(S_+ = S \cap P\) and \(S_- = S \cap N\). Then
\[
\sum_{x \in S} \chi(x) = |S_+| - |S_-|.
\]
Since \(0 \le |S_+| \le \lceil n/2 \rceil\) and
\(0 \le |S_-| \le \lfloor n/2 \rfloor\), we obtain
\[
-\left\lfloor \frac{n}{2} \right\rfloor
\le |S_+| - |S_-
\le \left\lceil \frac{n}{2} \right\rceil.
\]
Hence
\[
\left| \sum_{x \in S} \chi(x) \right|
\le \left\lceil \frac{n}{2} \right\rceil
\]
for every \(S \subseteq X\). Taking the maximum over
\(S \in \mathcal{P}(X)\) and then the minimum over colorings, we
conclude that
\[
\disc(\mathcal{P}(X)) \le \left\lceil \frac{n}{2} \right\rceil.
\]

Combining the lower and upper bounds yields
\[
\disc(\mathcal{P}(X)) = \left\lceil \frac{n}{2} \right\rceil.
\qquad \blacksquare
\]
\end{proof}

\end{document}
